\chapter{Introducción} % Este es el título del capítulo

La enseñanza de programación actualmente enfrenta un cambio de paradigma derivado del uso de herramientas de Inteligencia Artificial (IA) basadas en Modelos de Lenguaje de Gran Escala (LLMs), como ChatGPT, Claude o Gemini. Estas herramientas han demostrado ser muy eficientes para resolver tareas lógicas y algorítmicas, logrando un desempeño equivalente al de un estudiante de programación destacado.

No obstante, utilizar estas herramientas de manera incorrecta puede generar lo que algunas fuentes definen como la “paradoja cognitiva de la inteligencia artificial” \cite{Jose2025Cognitive}, un concepto que plantea que la IA puede funcionar como una herramienta que mejore la capacidad lógica e intelectual o como un factor de degradación cognitiva. En el ámbito educativo, este fenómeno se denomina como “descarga cognitiva” \cite{Jose2025Cognitive}, donde el estudiante delega completamente la resolución del problema, reduciendo la comprensión del problema y favoreciendo un aprendizaje superficial.

Estudios recientes sugieren que la dependencia de asistentes de IA puede generar un impacto negativo de los usuarios durante evaluaciones técnicas, afectando directamente en habilidades como la lectura y depuración de código \cite{Shen2026Skill}. Esto representa un riesgo estructural para la formación en ingeniería, donde desarrollar pensamiento lógico y la capacidad de resolver problemas son habilidades fundamentales. 

Este problema es un desafío para los docentes de ingeniería, ya que deben validar que el código entregado por el alumno refleja correctamente la lógica del problema y analizar si el código es generado por Inteligencia Artificial. La revisión manual es el estándar más confiable para evaluar el razonamiento detrás del código, pero en grupos numerosos o profesores con múltiples grupos asignados, esta estrategia es ineficiente y poco escalable \cite{Mekterovic2023Scaling}. Como consecuencia, algunos profesores recurren a herramientas de evaluación automática que priorizan la funcionalidad del programa sin analizar la estructura lógica interna de la solución.

Las herramientas tradicionales de detección de similitud, como MOSS o JPlag, presentan limitaciones al basarse principalmente en comparaciones textuales. Al no analizar la equivalencia semántica ni el flujo lógico del programa, se dificulta identificar soluciones que conservan la misma lógica algorítmica, pero presentan diferencias sintácticas. Además, los detectores de contenido generado por IA tienen resultados con una precisión variable en código, puesto que están diseñados para ser utilizados en texto natural \cite{WeberWulff2023Testing}. 

Debido a estas limitaciones, el problema técnico se centra en la ausencia de herramientas que permitan evaluar la equivalencia lógica del código y la identificación de patrones asociados al uso de IA generativa de manera simultánea. Frente a este escenario, el presente trabajo terminal tiene como objetivo desarrollar Graphito, un sistema de comparación de código en C/C++ mediante similitud semántica y análisis estilométrico para cursos introductorios de programación. El sistema busca funcionar como una herramienta de apoyo a la decisión docente, que permita tener evaluaciones más integrales.

Para evaluar la equivalencia semántica, se plantea el uso de un modelo de transformer preentrenado que se especializa en código. Se seleccionó por su capacidad de incorporar información estructural del código mediante grafos de flujos de datos, permitiendo capturar relaciones entre las variables internas del programa. Con este enfoque se facilita la detección de la similitud lógica más allá de las diferencias superficiales en la sintaxis o el estilo \cite{Guo2021GraphCodeBERT, MartinezGil2026Improving}.

Para complementar el análisis semántico, se plantea integrar un módulo que realice un análisis estilometrico mediante una red neuronal convolucional a nivel de caracteres (CharCNN), entrenada para distinguir patrones estilísticos asociados a código humano y código generado por modelos de lenguaje. Este componente es capaz de extraer patrones estructurales y de formato, obteniendo una “huella digital” del código obtenido mediante modelos de lenguaje \cite{Zhang2015CharCNN, Tihanyi2025DNA}. 

Mediante este enfoque de doble canal (semántico y estilométrico), Graphito se posiciona como una herramienta de apoyo a la evaluación académica en cursos introductorios de programación, contribuyendo como herramienta de apoyo para fortalecer la integridad académica y el desarrollo de habilidades cognitivas fundamentales en programación.


\section{Descripción del Problema}

El aprendizaje en programación representa uno de los principales desafíos en carreras referentes a las ciencias de la computación, principalmente en cursos introductorios de lenguajes como C y C++. Múltiples estudios señalan que los estudiantes principiantes enfrentan dificultades para desarrollar habilidades como la estructuración lógica de programas, el manejo de entradas invalidas y la consideración de casos limite. Por ejemplo, “Rainfall” es un ejercicio ampliamente utilizado para evaluar la comprensión algorítmica básica. Este ejercicio presenta una tasa de resolución correcta entre 45\% y 72\%, mientras que errores fundamentales como la división por cero aparecen en mas del 23\% de soluciones \cite{Seppala2015Rainfall}. Esta situación se agrava en contextos universitarios con grandes grupos, donde la evaluación manual de código no es nada escalable. 

Frente a este escenario, algunas instituciones educativas adoptaron el uso de sistemas de evaluación automática de código para gestionar la carga docente. Sin embargo, existe evidencia de que estos sistemas presentan grandes limitaciones. Un reciente análisis encontró que aproximadamente el 81\% de las herramientas de evaluación de código son completamente automatizados, pero el 66\% se centran en verificar la funcionalidad correcta del código mediante pruebas unitarias dinámicas. Este enfoque ignora aspectos fundamentales como el diseño, la legibilidad o la estructura lógica del programa \cite{Messer2024Automated}. Además, el 43\% de los sistemas analizados solo otorgan retroalimentación binaria (correcto / incorrecto), limitando la capacidad de los estudiantes para comprender sus errores y desarrollar habilidades con la corrección de los mismos \cite{Rocha2022Feedback}. 

La aparición de los Modelo de Lenguaje Grandes (LLM) y la Inteligencia Artificial Generativa (GenAI) han generado nuevas vulnerabilidades en el sistema educativo y han mostrado más vulnerabilidades en herramientas tradicionales. Estas herramientas de detección de plagio como Moss o JPlag, se basan principalmente la similitud textual o léxica entre programas. El problema rádica en que el código generado por IA es instrinsecamente estocastico y puede producir resultados funcionalmente equivalentes pero estructuralmente diferentes, dificultando la detección mediante métodos convencionales \cite{ACM2024Detecting}. 

En este contexto, estudios recientes han demostrado que el código obtenido mediante GenAI presenta una similitud promedio de apenas 0.56 respecto a soluciones existentes, permitiendo evadir los sistemas de detección basados en comparación textual \cite{Dickey2025Failure}. Como consecuencia, se estima que la tasa potencial de plagio en cursos de programación podría incrementar desde métricas tradicionales entre 10\% - 15\% hasta estimaciones del 75\% cuando los estudiantes usan herramientas de generación de código mediante Inteligencia Artificial \cite{Dickey2025Failure}.

Además de los problemas relacionados con la integridad académica, la dependencia excesiva de herramientas de generación de código mediante IA plantea riesgos muy importantes para el proceso de aprendizaje. Una investigación reciente muestra que los estudiantes que delegan excesivamente tareas de programación a sistemas de IA, obtuvieron en promedio 17\% menos rendimiento en pruebas de evaluación técnica, con tamaños de efecto estadísticamente significativo ($d = 0.738, p = 0.010$). Estos resultados indican una disminución real en la comprensión de los conceptos fundamentales de programación \cite{Shen2026Skill}.

La problemática se agrava en el contexto de universidades publicas mexicanas. Existen estudios sobre la normatividad institucional que revelan que, aunque muchas instituciones invierten en software de detección de plagio, existe una falta de marcos regulatorios claros. En un análisis de 33 instituciones de educación superior publicas en México, se encontró que solo 10 de ellas mencionan el plagio en su normatividad institucional y solo 4 cuentan con una definición formal del mismo. Esto demuestra una brecha entre las políticas de integridad académica declaradas y su implementación practica \cite{Morales2021Integridad}. Por otro lado, los factores socio-tecnologicos también influyen en el contexto educativo en México, ya que, aproximadamente 39.1\% de los estudiantes carece de acceso adecuado a internet, y cerca del 57.3\% cuenta con dificultades significativas para seguir instrucciones en entornos digitales de aprendizaje, incrementando la dependencia de soluciones tecnológicas automatizadas \cite{Zapata2021Mexico}.

En conjunto, estos factores evidencian una transformación en el entorno educativo en programación. Por un lado, los cursos introductorios enfrentan problemas como la falta de escalabilidad en la evaluación manual; mientras que, las herramientas de evaluacioni automática presentan limitaciones pedagógicas y técnicas importantes. Asimismo, el reciente auge de la Inteligencia Artificial Generativa ha introducido nuevas formas de producir código que desafían los métodos tradicionales de detección de plagio y plantean riesgos importantes para el desarrollo cognitivo de los estudiantes. 

Por lo tanto, surge la necesidad de desarrollas nuevas herramientas de análisis y comparación de programas más haya de la similitud textual, que incorpore representaciones profundas del código, usando estructuras sintácticas y grafos de flujo de datos, que permitan identificar similitudes semánticas y patrones de autoría en programas en C/C++. Debido a la ausencia de estas herramientas capaces de abordar el análisis desde ambos puntos de vista constituye el principal problema de investigación que se busca abordar en el presente trabajo terminal.

\section{Justificación}
El creciente uso de herramientas de Inteligencia Artificial Generativa en el contexto educativo ha generado un cambio en la manera en la que los estudiantes realizan tareas académicas complejas, especialmente en campos de programación. A pesar de que estas herramientas prometen mejorar la productividad y facilitan la resolución de problemas, estudios recientes indican que su empleo indiscriminado puede traer consecuencias negativas en los procesos de aprendizaje, afectando mayormente a alumnos que se encuentran en las primeras etapas de su formación profesional; impactando el desarrollo de habilidades como el razonamiento lógico, la resolución de problemas y la comprensión profunda de conceptos. En este contexto, es necesario analizar el impacto educativo de las herramientas de IA y desarrollar sistemas con el objetivo de mitigar los efectos negativos derivados de su uso indiscriminado.

El aprendizaje en disciplinas técnicas como la programación no depende únicamente de adquirir información, sino de la construcción progresiva de esquemas cognitivos complejos almacenados en la memoria a largo plazo \cite{lodge2026ai}, los cuales permiten el desarrollo de habilidades como el reconocimiento de patrones, el diagnóstico de errores y la resolución eficiente de problemas. Sin embargo, la aparición de herramientas de IA generativa presenta un riesgo en este proceso formativo, especialmente en entornos académicos donde su uso no se encuentra regulado.

Uno de los fenómenos más importantes es el descrito recientemente en \cite{lodge2026ai} como la descarga cognitiva perjudicial (\textit{cognitive offloading}), en la cual los estudiantes delegan a sistemas basados en LLMs las tareas cognitivas que son fundamentales para su proceso de aprendizaje. Cuando un alumno principiante solicita la solución completa de un problema de programación a una herramienta de IA, está evitando el trabajo cognitivo intrínseco, como la generación de ideas, la recuperación de conocimientos previos, el análisis del problema y la elaboración de la solución. Este proceso mental es el que le permite construir los esquemas de conocimiento necesarios para el dominio técnico de la disciplina. Dickey menciona en \cite{Dickey2025Failure} que al evadir este proceso mental, el estudiante experimenta una gratificación inmediata por terminar la tarea, pero, en consecuencia, atrofia su desarrollo cognitivo, lo que impide la formación de la base de conocimientos necesaria para el pensamiento crítico, la resolución de problemas futuros y la capacidad de generar soluciones creativas.

Las consecuencias de este fenómeno ya pudieron ser evidenciadas por estudios como el realizado por Barshay en 2024. En \cite{Jose2025Cognitive} se reportó un fenómeno denominado ``Paradoja del Rendimiento'', donde los estudiantes utilizaron herramientas de IA para practicar ejercicios y lograron contestar un 48\% más de problemas correctamente durante la práctica; a pesar de esto, obtuvieron un 17\% menos de calificación en pruebas de comprensión conceptual. Este resultado sugiere que el uso de la IA puede mejorar el rendimiento de los usuarios en tareas específicas a corto plazo, pero a la vez deteriorar la comprensión profunda del contenido.

Investigaciones realizadas en 2026 han reportado resultados similares mediante estudios controlados sobre el aprendizaje en programación. En \cite{Shen2026Skill} se muestra que los desarrolladores que utilizaron asistencia de IA para aprender nuevas bibliotecas de programación presentaron una reducción del 17\% en sus evaluaciones de conocimientos. Este efecto presentó una significancia estadística ($d$ de Cohen = 0.738, $p = 0.10$), mientras que en pruebas piloto el impacto fue considerablemente mayor, alcanzando ($d$ de Cohen = 1.7, $p = 0.003$) \cite{Shen2026Skill}, lo que indica que la dependencia temprana de herramientas de IA puede afectar directamente la adquisición de nuevas habilidades técnicas.

Asimismo, el investigador Kolhatkar en 2025 reportó que delegar tareas cognitivas a la IA redujo la capacidad de resolución de problemas independientes de los estudiantes en un 35\% \cite{lodge2026ai}, mientras que otros estudios documentaron que los alumnos que externalizan completamente las tareas de generación de código a modelos de IA obtienen puntajes extremadamente bajos en pruebas de comprensión, entre el 24\% y el 39\% \cite{Shen2026Skill}. Sumado a lo anterior, se ha demostrado que existe una relación directa entre la cantidad de código generado por modelos de lenguaje que un estudiante utiliza y la pérdida de aprendizaje en evaluaciones posteriores \cite{Dickey2025Failure}.

Aunado a lo anterior, se ha demostrado que el uso de estas herramientas es cada vez más generalizado. Estudios recientes reportan que al menos el 75\% de los estudiantes de segundo año de informática ya utilizan herramientas de generación de código, mientras que aproximadamente el 80\% de los universitarios confirman el empleo de estas tecnologías en su proceso de estudio \cite{Dickey2025Failure, lodge2026ai}. Teniendo conocimiento de las consecuencias negativas que conlleva la utilización indiscriminada de sistemas basados en IA, esta información resulta muy preocupante, ya que en entornos académicos puede generarse una ilusión de competencia que oculta deficiencias reales en la comprensión de conceptos. En este contexto, la falta de mecanismos institucionales capaces de responder a la transición tecnológica actual puede afectar la integridad académica, la equidad en la evaluación y la credibilidad del sistema educativo.

En el contexto de las universidades públicas mexicanas, ignorar el impacto de la IA generativa en la evaluación y la ausencia de mecanismos efectivos de supervisión puede generar lo que \cite{Morales2021Integridad} describe como un entorno de simulación académica, donde las malas prácticas se normalizan y se pasan por alto. Esta situación, sumada al ``profundo vacío normativo'' que existe en las universidades públicas en México —donde, de 33 instituciones, solo 10 mencionan explícitamente el término ``plagio'' en sus reglamentos y solo 4 tienen una definición clara del mismo—, evidencia la falta de un marco regulatorio para enfrentar nuevas formas de fraude académico \cite{Morales2021Integridad}.

Esta situación se agrava debido a que las instituciones actualmente carecen de estrategias contra el plagio asistido por inteligencia artificial, ya que muchas universidades optan por evitar el problema por limitaciones técnicas o para evitar conflictos institucionales \cite{Morales2021Integridad}. Sin embargo, esta falta de acción tiene consecuencias acumulativas, puesto que termina socavando ``el esfuerzo de los estudiantes honestos, distorsionando la justicia en las calificaciones, erosionando la credibilidad de la institución y devaluando el título profesional que expide'' \cite{Dickey2025Failure}.

La forma más intuitiva de abordar este problema es desde la perspectiva pedagógica: realizar revisiones manuales de código y efectuar entrevistas de defensa son los métodos más efectivos para validar la comprensión de la lógica y la autoría del alumno \cite{Dickey2025Failure}. Sin embargo, esta solución es logísticamente inviable para grupos masivos, ya que aplicar dichos métodos durante apenas 10 minutos para un grupo de 100 estudiantes exige más de 16 horas de trabajo para el docente, lo que representa un cuello de botella insostenible en la evaluación. Esta falta de escalabilidad genera una brecha de injusticia en la calificación, pues la revisión manual es subjetiva, propensa a la fatiga y a los sesgos del evaluador; el plagio puede pasar desapercibido, o dos profesores pueden tomar decisiones distintas ante una misma situación \cite{Dickey2025Failure}.

Ante este contexto, una posible solución reside en las herramientas tradicionales de detección de plagio de código como MOSS o JPlag, ya que permitirían automatizar el proceso para determinar si el alumno escribió el código por sí mismo. Sin embargo, estas herramientas fueron diseñadas en un contexto previo a la aparición de la IA generativa. Su funcionamiento se basa en la comparación de secuencias de tokens o patrones superficiales de código frente a repositorios conocidos, lo cual resulta poco eficiente cuando se aplica a código generado por modelos de inteligencia artificial \cite{saglam2024detecting}.

El código generado por inteligencia artificial es, por definición, un texto original que no existe en una base de datos previa, por lo que los métodos tradicionales son ciegos ante él \cite{Dickey2025Failure}. Como consecuencia, estas herramientas proveen una falsa seguridad: por un lado, cometen muchos falsos positivos (93\% de similitud) donde las soluciones correctas son naturalmente similares; por otro, los cambios superficiales como el reordenamiento de instrucciones o la inserción de código ``muerto'' evaden fácilmente al detector \cite{saglam2024detecting}. De hecho, estudios han demostrado que la eficacia de los detectores tradicionales puede caer por debajo del 25\%, e incluso por debajo del 10\% en escenarios reales. Esto evidencia que dichas herramientas están obsoletas frente a los sistemas de IA capaces de generar código \cite{saglam2024detecting}.

Frente a este escenario, surge la necesidad de adoptar nuevas estrategias que permitan cubrir dos necesidades fundamentales: garantizar el rigor académico y optimizar el tiempo del docente. Con este propósito, los Sistemas de Apoyo a la Decisión (\textit{Decision Support System}, DSS) representan una alternativa viable, ya que no buscan reemplazar el trabajo docente, sino automatizar el análisis de grandes volúmenes de información académica, permitiéndole al profesor identificar patrones sospechosos que requieran revisión manual. Diversos estudios reportan que la incorporación de estos sistemas puede reducir el tiempo de evaluación docente hasta en un 70\% \cite{technicalreview2025dss}, permitiendo que el profesor se enfoque en actividades pedagógicas de mayor valor, resolviendo el cuello de botella técnico y garantizando una evaluación justa, rigurosa y escalable.

En conjunto, la evidencia presentada demuestra que el uso indiscriminado de herramientas de inteligencia artificial generativa no solo constituye un problema pedagógico, sino también un problema institucional y tecnológico que afecta significativamente los procesos de evaluación en cursos introductorios de programación. Por un lado, la delegación de procesos cognitivos pone en riesgo la formación profesional de los estudiantes; por otro, las instituciones universitarias carecen de marcos regulatorios y herramientas tecnológicas adecuadas para afrontar el nuevo paradigma que presenta la inteligencia artificial en la educación actual. Aunado a esto, los sistemas tradicionales de detección de plagio están quedando obsoletos debido a su incapacidad para detectar código generado por modelos de lenguaje.

Por ello, se vuelve necesario explorar enfoques computacionales más avanzados que permitan analizar el código desde perspectivas más profundas, considerando su estructura lógica y los patrones de estilometría del autor. Con este enfoque, el presente trabajo propone el desarrollo de un sistema capaz de apoyar la evaluación académica mediante el análisis de código fuente utilizando algoritmos de inteligencia artificial, con el objetivo de preservar la integridad académica y fomentar el fortalecimiento de los procesos de aprendizaje de los estudiantes en el contexto actual.

\section{Motivación}

La motivación detrás de este trabajo terminal emana de una dualidad: la fascinación por las fronteras actuales de la Inteligencia Artificial y la preocupación ética por la formación de las nuevas generaciones de ingenieros. 

Desde una perspectiva técnica, nos impulsó el deseo de aplicar conceptos avanzados de Procesamiento de Lenguaje Natural (NLP) como \textit{embeddings} y mecanismos de atención en un dominio no convencional: el código fuente. La posibilidad de tratar un programa no solo como texto, sino como una estructura con semántica profunda, nos permitió explorar cómo las Redes Neuronales Convolucionales (CNN) pueden "ver" patrones estilísticos que escapan al ojo humano. En este sentido, el trabajo del investigador Martínez Gil \cite{Guo2021GraphCodeBERT} ha sido una fuente de inspiración fundamental, pues su uso innovador de \textit{GraphCodeBERT} nos demostró que es posible construir puentes entre la lógica algorítmica y el aprendizaje profundo para resolver problemas complejos de autoría.

En el ámbito personal y académico, este proyecto nace de una reflexión profunda sobre nuestra propia trayectoria en la ESCOM. Al inicio de nuestra formación, experimentamos el surgimiento de herramientas como ChatGPT, las cuales percibimos inicialmente como tutores ideales. Sin embargo, con el tiempo identificamos un riesgo latente: la tentación de recurrir al "camino fácil", resolviendo tareas de forma automatizada sin comprender la lógica subyacente. Esta dependencia genera un vacío de conocimiento que observamos con preocupación en el entorno educativo, donde la facilidad de implementación está sustituyendo peligrosamente a la capacidad de abstracción.

Consideramos que el futuro de la ingeniería de software no reside en memorizar sintaxis, sino en la capacidad de guiar a la IA mediante conceptos sólidos de algoritmos y estructuras de datos. Relacionamos esta visión con la labor de un arquitecto: aunque él no coloca cada ladrillo manualmente, debe poseer un dominio absoluto de la ingeniería estructural para dirigir la obra y garantizar que los cimientos sean sólidos. Sin ese conocimiento base, el "arquitecto de software" pierde su capacidad de supervisión y el enfoque del \textit{Human-in-the-loop} desaparece.

Por lo tanto, nuestra motivación principal es proporcionar a los docentes de la ESCOM una herramienta que no actúe como un simple sensor de plagio, sino como un instrumento de diagnóstico. Queremos que \textbf{Graphito} ayude a identificar el uso inadecuado de la IA en etapas críticas de aprendizaje, fomentando que los estudiantes recuperen el interés por construir su propia lógica y comprendan que, para manejar las herramientas del futuro, primero deben comprender los fundamentos que las sostienen.

\section{Objetivos}
\subsection{Objetivo General}
Implementar un sistema de apoyo a la evaluación docente en cursos introductorios de programación en C/C++ basado en comparacion semántica de código, que permita evaluar la similitud lógica entre programas y detectar posibles patrones de generación mediante inteligencia artificial, con el fin de optimizar el proceso de evaluación. 

\subsection{Objetivos específicos}
\begin{itemize}
    \item Seleccionar los modelos de lenguaje y arquitecturas de red neuronal adecuados mediante una evaluación comparativa del estado del arte, para fundamentar la viabilidad tecnológica de los módulos semántico y estilométrico.
    \item Diseñar la arquitectura de software del sistema mediante el modelado de diagramas estándar (UML/BPMN), definiendo la orquestación entre el cliente web, la API de servicios y la persistencia de datos.
    \item Desarrollar el módulo de análisis semántico integrando el modelo pre-entrenado GraphCodeBERT, automatizando la extracción de Grafos de Flujo de Datos (DFG) y el cálculo de la similitud del coseno entre representaciones vectoriales.
    \item Entrenar un modelo de clasificación estilométrica (CharCNN), conformando previamente un conjunto de datos (código humano vs. generado por IA), para alcanzar una precisión superior al 85\% en la detección de posibles patrones de inteligencia artificial generativa.
    \item Integrar los motores de inteligencia artificial en una plataforma web funcional, permitiendo al docente configurar parámetros de evaluación y visualizar gráficamente los índices de similitud.
    \item Evaluar el desempeño del sistema en un entorno de pruebas controlado, ejecutando simulaciones con un repositorio estático de ejercicios de programación, para medir su precisión técnica y tiempos de respuesta previo a un escenario de producción.
\end{itemize}